\chapter{الدراسة النظرية}

\begin{chaptersummary}
يستعرض هذا الفصل الأسس النظرية التي يستند إليها المشروع: أنماط هندسة البرمجيات،
وبروتوكولات الاتصال في الوقت الفعلي، والتعرف التلقائي على الكلام، والنماذج اللغوية.
\end{chaptersummary}

\section{مقدمة}
[مقدمة الفصل.]

\section{أنماط هندسة البرمجيات}

\subsection{المعمارية النظيفة \en{(Clean Architecture)}}
[المحتوى من المسودة.]

\subsection{معمارية الخدمات المصغرة \en{(Microservices Architecture)}}
[المحتوى من المسودة.]

\section{بروتوكولات وبنية الاتصال في الوقت الفعلي}

\subsection{بروتوكول \en{WebSockets}}
[المحتوى من المسودة.]

\subsection{الاتصال عبر الويب في الوقت الفعلي \en{(WebRTC)}}
[المحتوى من المسودة.]

\subsection{وحدة التوجيه الانتقائي \en{(SFU)}}
[المحتوى من المسودة.]

\section{التعرف التلقائي على الكلام \en{(ASR)}}

\subsection{أوضاع المعالجة: \en{Batch} مقابل \en{Streaming}}
[المحتوى من المسودة.]

\section{النماذج اللغوية}

\subsection{بنية المحول \en{(Transformer Architecture)}}
[المحتوى من المسودة.]

\subsection{النماذج اللغوية الكبيرة \en{(LLMs)} والتوليد المعزز بالاسترجاع \en{(RAG)}}
[المحتوى من المسودة.]

\section{الخاتمة}
[خاتمة الفصل.]

\transition{بعد عرض الأسس النظرية، نستعرض في الفصل التالي الأعمال المشابهة
والمنصات القائمة، لتحديد الفجوة التي يعالجها هذا المشروع.}
